\documentclass[accentcolor=tud4b,colorbacktitle,inverttitle,landscape,german,presentation,t]{tudbeamer}
\mathcode`\,="013B%L�scht das Leerzeichen nach dem Komma in Dezimalzahlen
%\usepackage{multimedia}

%\usepackage{media9}

%\usepackage{movie15}


\usepackage{hyperref}
\usepackage{bbding}
\usepackage{units}
\usepackage{epstopdf}
\usepackage{subfigure}
\usepackage{caption}
\usepackage{ngerman}
\usepackage{color}
\usepackage{rotating}
\usepackage[latin1]{inputenc}   % deutsche Umlaute
\usepackage[T1]{fontenc}        % Darstellung von Schrift als echte T1 Fonts und nicht als Bilder
% \usepackage{multimedia}
% \usepackage[square]{natbib}
% \usepackage{bibgerm}
\usepackage{tikz}

\usetikzlibrary{shapes.geometric,shapes.arrows,decorations.pathmorphing}
\usetikzlibrary{matrix,chains,scopes,positioning,arrows,fit}
\usetikzlibrary{calc}
\usetikzlibrary{patterns}

% \usetikzlibrary{patterns}
 \usetikzlibrary{matrix}
 \usetikzlibrary{decorations.pathmorphing}
 \usetikzlibrary{arrows,shapes.geometric} % fadings, ???
 \usepackage{fancybox}
 \usepackage{amsmath}            % Mathematik-Umgebungen
\usepackage{bm}                 % fette Mathe Buchstaben
\usepackage{amssymb}            % spezielle Mathe-Symbole


\usepackage{color}

%
% Absolute positioning
% beamer paper size: 12.80cm x 9.60cm
%
\usepackage[overlay,absolute]{textpos}
\setlength{\TPHorizModule}{10mm}
\setlength{\TPVertModule}{\TPHorizModule}
\textblockorigin{5mm}{25.5mm}
\setlength{\parindent}{0pt}

\newcommand{\leftbox}[1]{
 \begin{textblock}{5.7}(0, 0)
   #1
 \end{textblock}
}

\newcommand{\rightbox}[1]{
 \begin{textblock}{5.7}(6.1, 0)
   #1
 \end{textblock}
}

\newcommand{\fullbox}[1]{
 \begin{textblock}{11.8}(0, 0)
   #1
 \end{textblock}
}

% \usepackage{calc}
\newenvironment{arbbox}[2]
{%\setlength{\boxwidth}{11.8cm minus \real{#1}}
\begin{textblock}{11.8cm}(#1, #2)}
{\end{textblock}}




%
% Macros
%
% \renewcommand{\vec}[1]{\ensuremath{\bm{#1}}}
\renewcommand{\vec}[1]{\ensuremath{\boldsymbol{#1}}}

\begin{document}
\title[]{Titel}
\subtitle{Namen}

\author[Gruppe X]{Gruppe X}
\institute[Regelungsmethoden und Robotik | Projektseminar Robotik]{Regelungsmethoden und Robotik | Projektseminar Robotik}




% \logo{\color{tudtextaccent}TU Darmstadt-}
\logo{\includegraphics{rmr2}}

\date{XX.YY.2016}
%Titelfolie%%%%%%%%%%%%%%%%%%%%%%%%%%%%%%%%%%%%%%%%%%%%%%%%%%%%%%%%%%%%%%%%%%%%%%%%%%%%%%%%%%%%%%%%%%%%%%%%%%%%%%%%%%%%%%%%%%%%%%%%%%%%%%%%%%%%%%%%%%%%%%%%%%%
\begin{titleframe}





\end{titleframe}

%Hauptfolien%%%%%%%%%%%%%%%%%%%%%%%%%%%%%%%%%%%%%%%%%%%%%%%%%%%%%%%%%%%%%%%%%%%%%%%%%%%%%%%%%%%%%%%%%%%%%%%%%%%%%%%%%%%%%%%%%%%%%%%%%%%%%%%%%%%%%%%%%%%%%%%%%%%
%%%%%%%%%%

\begin{frame}{\\�bersicht}
 \leftbox{
 \vspace{1.5cm}
 \begin{itemize}
  \item Something
  \item Something
  \item Something
  \item ...
 
 \end{itemize}
}

\end{frame}








\end{document}
